\begin{table*}[htbp]
\centering
\small
\caption{Conceptual definitions and terminology across the digital patient spectrum.}
\label{tab:concepts}
\begin{tabularx}{\textwidth}{@{}>{\raggedright\arraybackslash}p{0.20\textwidth}YY@{}}
\toprule
\textbf{Conceptual dimension} & \textbf{Current extracted signal} & \textbf{Interpretive use} \\
\midrule
Primary terminology & Digital twin (47); Virtual patient (5); Human digital twin (5); Patient digital twin (4); In silico patient / trial (3) & Shows the distribution of the first identifiable primary construct per review record. \\
Patient specificity & Patient-specific/individualized (56); generic/population/archetype (28). & Separates individualized models from population-level or archetypal constructs; categories may overlap. \\
Temporal orientation & Dynamic/continuous (54); explicitly static (15). & Tests whether digital twin terminology is accompanied by longitudinal updating. \\
Data coupling & Bidirectional (33); one-way (39). & Distinguishes model-, shadow-, and twin-like coupling claims. \\
Hierarchy & Hierarchy levels are converted to an ordinal cell-to-system-of-systems scale for Figure 2; multi-level entries use the highest coded level. & Makes the spectrum plot reproducible while preserving the row-level source text. \\
\bottomrule
\end{tabularx}
\end{table*}
